% ch3-systemes-lineaires.tex — Chapitre 3 : Systèmes d'équations linéaires et élimination de Gauss
% Ce chapitre développe la théorie des systèmes linéaires, l'algorithme de
% Gauss, le théorème de Rouché–Capelli, et la structure de l'ensemble des
% solutions (sous-espace affine). Des applications aux réseaux de flux et
% circuits électriques sont présentées.

\chapter{Systèmes d'équations linéaires et élimination de Gauss}\label{ch:systemes-lineaires}

\section*{Introduction}
\addcontentsline{toc}{section}{Introduction}

Les systèmes d'équations linéaires\index{système linéaire} apparaissent
dans pratiquement tous les domaines des mathématiques appliquées et des
sciences. En physique, l'analyse des circuits électriques conduit à des
systèmes linéaires via les lois de Kirchhoff\index{Kirchhoff (lois de)}~;
en économie, les modèles d'équilibre de Leontief\index{Leontief (modèle de)}
se ramènent à la résolution d'un système $A\vec{x}=\vec{b}$~; en
informatique, les moteurs de recherche, la compression d'images et
l'apprentissage automatique reposent tous, in fine, sur la résolution
efficace de systèmes linéaires de grande taille.

\medskip

Ce chapitre introduit un outil fondamental pour l'étude de ces systèmes~:
l'\emph{élimination de Gauss}\index{élimination de Gauss}, un algorithme
systématique qui ramène tout système linéaire à une forme échelonnée
facile à résoudre. Nous établirons le lien entre le rang d'une matrice et
l'existence de solutions (théorème de Rouché--Capelli), et nous montrerons
que l'ensemble des solutions d'un système linéaire possède une structure
géométrique précise~: c'est un sous-espace affine.

\medskip

\begin{center}
\begin{tikzpicture}[scale=1.1, >=Stealth]
  % Système 2x2 : intersection de deux droites
  \begin{scope}
    \draw[->] (-0.5,0) -- (4.2,0) node[right] {$x_1$};
    \draw[->] (0,-0.5) -- (0,3.7) node[above] {$x_2$};
    % Droite 1: x1 + x2 = 3  =>  x2 = 3 - x1
    \draw[thick, blue!70!black, domain=0:3.5] plot (\x, {3 - \x})
      node[right, font=\small] {$x_1+x_2=3$};
    % Droite 2: 2x1 - x2 = 0  =>  x2 = 2*x1
    \draw[thick, red!70!black, domain=0:1.8] plot (\x, {2*\x})
      node[right, font=\small] {$2x_1-x_2=0$};
    % Intersection
    \filldraw[black] (1,2) circle (2pt) node[above right, font=\small] {$(1,2)$};
    \node[below, font=\small] at (2,-0.8) {Système compatible déterminé};
  \end{scope}
  % Système 2x2 : droites parallèles (incompatible)
  \begin{scope}[xshift=6.5cm]
    \draw[->] (-0.5,0) -- (4.2,0) node[right] {$x_1$};
    \draw[->] (0,-0.5) -- (0,3.7) node[above] {$x_2$};
    % Droite 1: x1 + x2 = 2
    \draw[thick, blue!70!black, domain=0:2.5] plot (\x, {2 - \x})
      node[right, font=\small] {$x_1+x_2=2$};
    % Droite 2: x1 + x2 = 4
    \draw[thick, red!70!black, domain=0:3.8] plot (\x, {4 - \x})
      node[right, font=\small] {$x_1+x_2=4$};
    \node[below, font=\small] at (2,-0.8) {Système incompatible};
  \end{scope}
\end{tikzpicture}
\captionof{figure}{Interprétation géométrique d'un système de deux équations à deux inconnues~: solution unique (gauche) ou absence de solution (droite).}
\label{fig:systemes-2x2-intro}
\end{center}

% ============================================================
\section{Systèmes d'équations linéaires}\label{sec:systemes-def}
% ============================================================

\subsection{Définition et notation}\label{subsec:systemes-notation}

\begin{definition}{Système d'équations linéaires}{def:systeme-lineaire}
  Un \emph{système de $m$ équations linéaires à $n$ inconnues}\index{système linéaire!définition}
  sur un corps $\K$ est un ensemble d'équations de la forme~:
  \begin{equation}\label{eq:systeme-general}
    \left\{
    \begin{aligned}
      a_{11}x_1 + a_{12}x_2 + \cdots + a_{1n}x_n &= b_1,\\
      a_{21}x_1 + a_{22}x_2 + \cdots + a_{2n}x_n &= b_2,\\
      &\;\,\vdots\\
      a_{m1}x_1 + a_{m2}x_2 + \cdots + a_{mn}x_n &= b_m,
    \end{aligned}
    \right.
  \end{equation}
  où les \emph{coefficients} $a_{ij}\in\K$ et les \emph{seconds membres}
  $b_i\in\K$ sont donnés, et les $x_j$ sont les \emph{inconnues}.
\end{definition}

\subsection{Écriture matricielle}\label{subsec:ecriture-matricielle}

Le système~\eqref{eq:systeme-general} se réécrit de manière compacte
sous la forme
\begin{equation}\label{eq:Axb}
  A\vec{x}=\vec{b},
\end{equation}
où
\[
  A = \begin{pmatrix}
    a_{11} & a_{12} & \cdots & a_{1n}\\
    a_{21} & a_{22} & \cdots & a_{2n}\\
    \vdots & \vdots & \ddots & \vdots\\
    a_{m1} & a_{m2} & \cdots & a_{mn}
  \end{pmatrix}\in\Mat_{m,n}(\K),\qquad
  \vec{x} = \begin{pmatrix}x_1\\\vdots\\x_n\end{pmatrix}\in\K^n,\qquad
  \vec{b} = \begin{pmatrix}b_1\\\vdots\\b_m\end{pmatrix}\in\K^m.
\]
La matrice $A$ est la \emph{matrice des coefficients}\index{matrice!des coefficients}
du système. La \emph{matrice augmentée}\index{matrice!augmentée} est la
matrice $(A\mid\vec{b})\in\Mat_{m,n+1}(\K)$ obtenue en juxtaposant $\vec{b}$
comme colonne supplémentaire~:
\begin{equation}\label{eq:matrice-augmentee}
  (A\mid\vec{b}) = \left(\begin{array}{cccc|c}
    a_{11} & a_{12} & \cdots & a_{1n} & b_1\\
    a_{21} & a_{22} & \cdots & a_{2n} & b_2\\
    \vdots & \vdots & \ddots & \vdots & \vdots\\
    a_{m1} & a_{m2} & \cdots & a_{mn} & b_m
  \end{array}\right).
\end{equation}

\begin{definition}{Solution d'un système}{def:solution-systeme}
  Un vecteur $\vec{s}=(s_1,\ldots,s_n)\in\K^n$ est une \emph{solution}%
  \index{solution d'un système} du système~\eqref{eq:Axb} si $A\vec{s}=\vec{b}$.
  L'\emph{ensemble des solutions}\index{ensemble des solutions} est
  \[
    \cS \deff \setcond{\vec{x}\in\K^n}{A\vec{x}=\vec{b}}.
  \]
  Le système est dit \emph{compatible}\index{système linéaire!compatible}
  (ou \emph{consistant}) si $\cS\neq\varnothing$, et
  \emph{incompatible}\index{système linéaire!incompatible} sinon.
\end{definition}

\subsection{Systèmes homogènes et non homogènes}\label{subsec:homogene-non}

\begin{definition}{Système homogène}{def:systeme-homogene}
  Le système~\eqref{eq:Axb} est dit \emph{homogène}\index{système linéaire!homogène}
  si $\vec{b}=\vzero$, c'est-à-dire si toutes les constantes du second
  membre sont nulles. Le système homogène associé à~\eqref{eq:Axb} est~:
  \begin{equation}\label{eq:systeme-homogene}
    A\vec{x}=\vzero.
  \end{equation}
\end{definition}

\begin{remark}{Compatibilité d'un système homogène}{rem:homogene-compatible}
  Un système homogène est toujours compatible car $\vec{x}=\vzero$ est
  solution. Cette solution est dite \emph{triviale}\index{solution!triviale}.
  La question est donc de savoir s'il existe des solutions
  \emph{non triviales}.
\end{remark}

\begin{proposition}{Solutions d'un système homogène}{prop:sol-homogene-sev}
  L'ensemble des solutions du système homogène $A\vec{x}=\vzero$ est un
  sous-espace vectoriel\index{sous-espace vectoriel} de $\K^n$. C'est
  exactement le noyau $\Ker(A)$ de l'application linéaire
  $\vec{x}\mapsto A\vec{x}$.
\end{proposition}

\begin{proof}
  Posons $f\colon\K^n\to\K^m$, $\vec{x}\mapsto A\vec{x}$. L'application
  $f$ est linéaire (cf.~
ef{ch:applications-lineaires}). L'ensemble des
  solutions de $A\vec{x}=\vzero$ est $\Ker(f)$, qui est un sous-espace
  vectoriel de $\K^n$ d'après la 
ef{ch:applications-lineaires}.
\end{proof}

% ============================================================
\section{Opérations élémentaires sur les lignes}\label{sec:operations-elementaires}
% ============================================================

L'idée clé de l'élimination de Gauss est de transformer un système en un
système \emph{équivalent} (c'est-à-dire ayant le même ensemble de solutions)
mais de forme plus simple. Les transformations autorisées sont les
\emph{opérations élémentaires sur les lignes}\index{opérations élémentaires}.

\begin{definition}{Opérations élémentaires sur les lignes}{def:ops-elem}
  Soit $M\in\Mat_{m,n}(\K)$. On appelle \emph{opérations élémentaires
  sur les lignes}\index{opérations élémentaires!sur les lignes} les trois
  types de transformations suivants~:
  \begin{enumerate}[label=(\roman*)]
    \item \textbf{Échange de deux lignes~:}\index{échange de lignes}
      $L_i \leftrightarrow L_j$ \quad ($i\neq j$).
    \item \textbf{Multiplication d'une ligne par un scalaire non nul~:}%
      \index{multiplication d'une ligne}
      $L_i \leftarrow \lambda L_i$ \quad ($\lambda\in\K^{*}$).
    \item \textbf{Ajout d'un multiple d'une ligne à une autre~:}%
      \index{combinaison de lignes}
      $L_i \leftarrow L_i + \lambda L_j$ \quad ($\lambda\in\K$, $i\neq j$).
  \end{enumerate}
\end{definition}

\begin{proposition}{Conservation de l'ensemble des solutions}{prop:ops-elem-equiv}
  Les opérations élémentaires sur les lignes ne modifient pas l'ensemble
  des solutions du système linéaire. Autrement dit, si la matrice augmentée
  $(A'\mid\vec{b}')$ est obtenue à partir de $(A\mid\vec{b})$ par une suite
  d'opérations élémentaires, alors les systèmes $A\vec{x}=\vec{b}$ et
  $A'\vec{x}=\vec{b}'$ ont le même ensemble de solutions.
\end{proposition}

\begin{proof}
  Chaque opération élémentaire correspond à la multiplication à gauche par
  une matrice inversible\index{matrice!inversible} $E\in\GL_{m}(\K)$ (matrice
  élémentaire). Ainsi $A'\vec{x}=\vec{b}'$ s'écrit $EA\vec{x}=E\vec{b}$.
  Comme $E$ est inversible, $EA\vec{x}=E\vec{b}$ si et seulement si
  $A\vec{x}=\vec{b}$ (il suffit de multiplier à gauche par $\inv{E}$).
  Par composition, toute suite finie d'opérations élémentaires préserve
  l'ensemble des solutions.
\end{proof}

\begin{example}{Matrices élémentaires}{ex:matrices-elem}
  Dans $\Mat_{3,3}(\R)$, les trois types de matrices élémentaires sont~:
  \[
    E_{L_1\leftrightarrow L_2}=\begin{pmatrix}0&1&0\\1&0&0\\0&0&1\end{pmatrix},\qquad
    E_{L_2\leftarrow 3L_2}=\begin{pmatrix}1&0&0\\0&3&0\\0&0&1\end{pmatrix},\qquad
    E_{L_3\leftarrow L_3+2L_1}=\begin{pmatrix}1&0&0\\0&1&0\\2&0&1\end{pmatrix}.
  \]
  Chacune est inversible~; par exemple,
  $\inv{E}_{L_3\leftarrow L_3+2L_1}=E_{L_3\leftarrow L_3-2L_1}$.
\end{example}

% ============================================================
\section{Forme échelonnée et forme échelonnée réduite}\label{sec:echelon}
% ============================================================

\begin{definition}{Pivot}{def:pivot}
  Soit $M=(m_{ij})$ une matrice. Pour une ligne $L_i$ non nulle, le
  \emph{pivot}\index{pivot} de $L_i$ est le premier coefficient non nul de
  cette ligne, c'est-à-dire $m_{ij}$ où $j=\min\setcond{k}{m_{ik}\neq 0}$.
\end{definition}

\begin{definition}{Forme échelonnée (en lignes)}{def:forme-echelonnee}
  Une matrice est sous \emph{forme échelonnée}\index{forme échelonnée}
  (ou \emph{forme échelonnée en lignes}, \emph{row echelon form}) si~:
  \begin{enumerate}[label=(\roman*)]
    \item toutes les lignes nulles sont en bas de la matrice~;
    \item le pivot de chaque ligne non nulle est strictement à droite du
      pivot de la ligne précédente.
  \end{enumerate}
  Si, de plus~:
  \begin{enumerate}[label=(\roman*),resume]
    \item chaque pivot vaut~$1$~;
    \item chaque pivot est le seul élément non nul de sa colonne~;
  \end{enumerate}
  la matrice est sous \emph{forme échelonnée réduite}\index{forme échelonnée réduite}
  (\emph{reduced row echelon form}, abrégée RREF\index{RREF}).
\end{definition}

\begin{example}{Formes échelonnées}{ex:formes-echelonnees}
  La matrice suivante est sous forme échelonnée (les pivots sont encadrés)~:
  \[
    \begin{pmatrix}
      \boxed{2} & 3 & -1 & 7\\
      0 & \boxed{1} & 4 & 2\\
      0 & 0 & 0 & \boxed{5}\\
      0 & 0 & 0 & 0
    \end{pmatrix}.
  \]
  La forme échelonnée réduite correspondante est~:
  \[
    \begin{pmatrix}
      \boxed{1} & 0 & -13 & 0\\
      0 & \boxed{1} & 4 & 0\\
      0 & 0 & 0 & \boxed{1}\\
      0 & 0 & 0 & 0
    \end{pmatrix}.
  \]
\end{example}

\begin{theorem}{Unicité de la forme échelonnée réduite}{thm:unicite-rref}
  Toute matrice $A\in\Mat_{m,n}(\K)$ est équivalente par lignes à une
  unique matrice sous forme échelonnée réduite.
\end{theorem}

\begin{proof}
  \emph{Existence.} L'algorithme de Gauss--Jordan (
ef{sec:gauss-jordan})
  produit une forme échelonnée réduite en un nombre fini d'opérations
  élémentaires.

  \emph{Unicité.} Soient $R$ et $R'$ deux formes échelonnées réduites
  de~$A$, de sorte que $R=PA$ et $R'=QA$ pour des matrices inversibles $P$
  et $Q$. Comme $R$ et $R'$ sont équivalentes par lignes, il existe $S$
  inversible telle que $R'=SR$. Montrons que $R=R'$.

  Les colonnes pivots de $R$ et $R'$ sont identiques~: en effet, la colonne
  $j$ de $R$ (resp.~$R'$) est une colonne pivot si et seulement si elle
  n'est pas combinaison linéaire des colonnes précédentes de $R$
  (resp.~$R'$), propriété qui ne dépend que de l'espace des colonnes, lequel
  est le même pour $R$ et $R'$ (car elles sont équivalentes par lignes et
  l'espace des lignes est préservé, ce qui entraîne que les relations de
  dépendance entre colonnes sont identiques).

  Notons $j_1<j_2<\cdots<j_r$ les indices des colonnes pivots communs.
  Pour $k=1,\ldots,r$, la colonne $j_k$ de $R$ est le vecteur $\vece{k}$
  de $\R^m$ (vecteur avec $1$ en position $k$ et $0$ ailleurs), et de même
  pour $R'$. Pour une colonne non pivot $j$, les coefficients de la colonne
  $j$ dans $R$ sont déterminés de façon unique par les relations de
  dépendance linéaire entre les colonnes de $A$, qui sont les mêmes pour
  $R$ et $R'$. Donc $R=R'$.
\end{proof}

% ============================================================
\section{Élimination de Gauss}\label{sec:gauss}
% ============================================================

L'\emph{élimination de Gauss}\index{élimination de Gauss} (ou
\emph{méthode du pivot de Gauss}\index{pivot de Gauss}) est un algorithme
qui transforme la matrice augmentée $(A\mid\vec{b})$ en une forme
échelonnée par des opérations élémentaires sur les lignes.

\subsection{Algorithme}\label{subsec:algo-gauss}

\begin{enumerate}[label=\textbf{Étape \arabic*.},leftmargin=*]
  \item\label{step:gauss-pivot} \textbf{Choix du pivot.} On considère la
    colonne la plus à gauche qui contient un élément non nul sous la ligne
    courante (ou sur cette ligne). Si un tel élément se trouve sur une ligne
    $L_k$ avec $k>i$ (où $i$ est la ligne courante), on effectue l'échange
    $L_i\leftrightarrow L_k$.
  \item\label{step:gauss-elim} \textbf{Élimination.} Pour chaque ligne
    $L_j$ avec $j>i$, on effectue
    $L_j\leftarrow L_j - \frac{a_{ji}}{a_{ii}}L_i$,
    ce qui annule le coefficient en position $(j,i)$.
  \item\label{step:gauss-iter} \textbf{Itération.} On passe à la ligne
    suivante $i\leftarrow i+1$ et on recommence à l'étape~\ref{step:gauss-pivot}
    tant qu'il reste des lignes et des colonnes à traiter.
\end{enumerate}

Le résultat est une matrice sous forme échelonnée. On résout alors le
système par \emph{substitution arrière}\index{substitution arrière}
(\emph{back substitution}).

\subsection{Exemple détaillé}\label{subsec:ex-gauss}

\begin{example}{Élimination de Gauss sur un système $3\times 3$}{ex:gauss-3x3}
  Résolvons le système~:
  \[
    \left\{
    \begin{aligned}
      2x_1 + x_2 - x_3 &= 8,\\
      -3x_1 - x_2 + 2x_3 &= -11,\\
      -2x_1 + x_2 + 2x_3 &= -3.
    \end{aligned}
    \right.
  \]
  Matrice augmentée~:
  \[
    \left(\begin{array}{ccc|c}
      2 & 1 & -1 & 8\\
      -3 & -1 & 2 & -11\\
      -2 & 1 & 2 & -3
    \end{array}\right).
  \]

  \textbf{Étape 1~:} $L_2\leftarrow L_2+\frac{3}{2}L_1$,\quad
  $L_3\leftarrow L_3+L_1$~:
  \[
    \left(\begin{array}{ccc|c}
      2 & 1 & -1 & 8\\[3pt]
      0 & \frac{1}{2} & \frac{1}{2} & 1\\[3pt]
      0 & 2 & 1 & 5
    \end{array}\right).
  \]

  \textbf{Étape 2~:} $L_3\leftarrow L_3-4L_2$~:
  \[
    \left(\begin{array}{ccc|c}
      2 & 1 & -1 & 8\\[3pt]
      0 & \frac{1}{2} & \frac{1}{2} & 1\\[3pt]
      0 & 0 & -1 & 1
    \end{array}\right).
  \]

  \textbf{Substitution arrière~:}
  \begin{itemize}
    \item $L_3$~: $-x_3=1$, d'où $x_3=-1$.
    \item $L_2$~: $\frac{1}{2}x_2+\frac{1}{2}(-1)=1$, d'où $x_2=3$.
    \item $L_1$~: $2x_1+3-(-1)=8$, d'où $x_1=2$.
  \end{itemize}
  L'unique solution est $(x_1,x_2,x_3)=(2,3,-1)$.
\end{example}

\begin{example}{Système avec infinité de solutions}{ex:gauss-infini}
  Résolvons~:
  \[
    \left\{
    \begin{aligned}
      x_1 - 2x_2 + x_3 &= 0,\\
      2x_1 + x_2 - x_3 &= 3,\\
      4x_1 - 3x_2 + x_3 &= 3.
    \end{aligned}
    \right.
  \]
  Matrice augmentée~:
  \[
    \left(\begin{array}{ccc|c}
      1 & -2 & 1 & 0\\
      2 & 1 & -1 & 3\\
      4 & -3 & 1 & 3
    \end{array}\right)
    \xrightarrow{L_2\leftarrow L_2-2L_1}
    \left(\begin{array}{ccc|c}
      1 & -2 & 1 & 0\\
      0 & 5 & -3 & 3\\
      4 & -3 & 1 & 3
    \end{array}\right)
    \xrightarrow{L_3\leftarrow L_3-4L_1}
    \left(\begin{array}{ccc|c}
      1 & -2 & 1 & 0\\
      0 & 5 & -3 & 3\\
      0 & 5 & -3 & 3
    \end{array}\right).
  \]
  Puis $L_3\leftarrow L_3-L_2$~:
  \[
    \left(\begin{array}{ccc|c}
      1 & -2 & 1 & 0\\
      0 & 5 & -3 & 3\\
      0 & 0 & 0 & 0
    \end{array}\right).
  \]
  On a deux pivots et trois inconnues~: $x_3$ est une variable libre.
  En posant $x_3=t\in\R$~:
  \[
    x_2=\frac{3+3t}{5},\qquad
    x_1=2\cdot\frac{3+3t}{5}-t=\frac{6+6t-5t}{5}=\frac{6+t}{5}.
  \]
  L'ensemble des solutions est~:
  \[
    \cS=\setcond{\frac{1}{5}\begin{pmatrix}6\\3\\0\end{pmatrix}
      +t\,\frac{1}{5}\begin{pmatrix}1\\3\\5\end{pmatrix}}{t\in\R}
    =\set{\vec{x}_p+t\,\vec{v}\mid t\in\R},
  \]
  qui est une droite affine de $\R^3$.
\end{example}

% ============================================================
\section{Élimination de Gauss--Jordan}\label{sec:gauss-jordan}
% ============================================================

L'\emph{élimination de Gauss--Jordan}\index{élimination de Gauss--Jordan}
prolonge l'élimination de Gauss pour obtenir la forme échelonnée réduite
(RREF). Après avoir obtenu une forme échelonnée par l'algorithme de Gauss,
on effectue les étapes supplémentaires suivantes~:

\begin{enumerate}[label=\textbf{Étape \arabic*.},leftmargin=*]
  \item\label{step:gj-norm} \textbf{Normalisation des pivots.} Pour chaque
    ligne non nulle $L_i$ de pivot $a_{ij}$, on effectue
    $L_i\leftarrow\frac{1}{a_{ij}}L_i$ pour rendre le pivot égal à~$1$.
  \item\label{step:gj-elim-haut} \textbf{Élimination vers le haut.}
    En partant du dernier pivot et en remontant, pour chaque pivot en
    position $(i,j)$, on annule tous les coefficients au-dessus du pivot~:
    $L_k\leftarrow L_k-a_{kj}L_i$ pour $k<i$.
\end{enumerate}

\begin{example}{Gauss--Jordan complet}{ex:gauss-jordan}
  Reprenons la matrice échelonnée de l'
ef{ex:ex:gauss-3x3}~:
  \[
    \left(\begin{array}{ccc|c}
      2 & 1 & -1 & 8\\[3pt]
      0 & \frac{1}{2} & \frac{1}{2} & 1\\[3pt]
      0 & 0 & -1 & 1
    \end{array}\right).
  \]
  \textbf{Normalisation~:}
  $L_1\leftarrow\frac{1}{2}L_1$, $L_2\leftarrow 2L_2$,
  $L_3\leftarrow -L_3$~:
  \[
    \left(\begin{array}{ccc|c}
      1 & \frac{1}{2} & -\frac{1}{2} & 4\\[3pt]
      0 & 1 & 1 & 2\\[3pt]
      0 & 0 & 1 & -1
    \end{array}\right).
  \]
  \textbf{Élimination vers le haut~:}
  $L_2\leftarrow L_2-L_3$, puis $L_1\leftarrow L_1+\frac{1}{2}L_3$,
  puis $L_1\leftarrow L_1-\frac{1}{2}L_2$~:
  \[
    \left(\begin{array}{ccc|c}
      1 & 0 & 0 & 2\\
      0 & 1 & 0 & 3\\
      0 & 0 & 1 & -1
    \end{array}\right).
  \]
  On lit directement la solution~: $(x_1,x_2,x_3)=(2,3,-1)$.
\end{example}

% ============================================================
\section{Rang d'une matrice}\label{sec:rang-ch3}
% ============================================================

\begin{definition}{Rang d'une matrice}{def:rang-ch3}
  Le \emph{rang}\index{rang} d'une matrice $A\in\Mat_{m,n}(\K)$, noté
  $\rg(A)$, est le nombre de pivots dans toute forme échelonnée de $A$.
  De manière équivalente~:
  \begin{enumerate}[label=(\roman*)]
    \item $\rg(A)$ est la dimension de l'espace vectoriel engendré par les
      lignes de $A$ (rang ligne)\index{rang!ligne}~;
    \item $\rg(A)$ est la dimension de l'espace vectoriel engendré par les
      colonnes de $A$ (rang colonne)\index{rang!colonne}~;
    \item $\rg(A)=\dim(\im(f))$, où $f\colon\K^n\to\K^m$,
      $\vec{x}\mapsto A\vec{x}$.
  \end{enumerate}
\end{definition}

\begin{remark}{Rang ligne = rang colonne}{rem:rang-egal}
  L'égalité entre le rang ligne et le rang colonne est un résultat
  fondamental. On l'établit en remarquant que les opérations élémentaires
  sur les lignes ne changent pas les relations de dépendance linéaire entre
  les colonnes, et que dans une forme échelonnée, les colonnes pivots
  forment une famille libre de cardinal égal au nombre de lignes non nulles.
\end{remark}

\begin{proposition}{Propriétés du rang}{prop:proprietes-rang}
  Soit $A\in\Mat_{m,n}(\K)$. Alors~:
  \begin{enumerate}[label=(\roman*)]
    \item $0\leqslant\rg(A)\leqslant\min(m,n)$.
    \item $\rg(A)=\rg(\transpose{A})$.
    \item Si $P\in\GL_m(\K)$ et $Q\in\GL_n(\K)$, alors $\rg(PAQ)=\rg(A)$.
    \item Pour $A$ carrée d'ordre $n$~: $A$ est inversible si et seulement
      si $\rg(A)=n$.
  \end{enumerate}
\end{proposition}

\begin{proof}
  (i) est immédiat. (ii) découle de l'égalité rang ligne = rang colonne.
  (iii)~: la multiplication à gauche par $P$ inversible correspond à des
  opérations élémentaires sur les lignes (qui ne changent pas le rang), et
  la multiplication à droite par $Q$ inversible correspond à des opérations
  élémentaires sur les colonnes (qui ne changent pas le rang non plus).
  (iv)~: $A$ est inversible si et seulement si sa forme échelonnée réduite
  est $I_n$, c'est-à-dire si et seulement si elle a $n$ pivots.
\end{proof}

% ============================================================
\section{Théorème de Rouché--Capelli}\label{sec:rouche-capelli}
% ============================================================

Le théorème suivant, connu sous le nom de \emph{théorème de
Rouché--Capelli}\index{Rouché--Capelli (théorème de)} (ou
\emph{Kronecker--Capelli}\index{Kronecker--Capelli (théorème de)} dans
la tradition allemande), fournit le critère fondamental d'existence de
solutions.

\begin{theorem}{Théorème de Rouché--Capelli}{thm:rouche-capelli}
  Soit $A\vec{x}=\vec{b}$ un système linéaire avec $A\in\Mat_{m,n}(\K)$
  et $\vec{b}\in\K^m$. Alors~:
  \begin{enumerate}[label=(\roman*)]
    \item Le système est compatible si et seulement si
      $\rg(A)=\rg(A\mid\vec{b})$.
    \item Lorsque le système est compatible, la dimension de l'ensemble des
      solutions du système homogène associé $A\vec{x}=\vzero$ est
      $n-\rg(A)$.
    \item En particulier, la solution est unique si et seulement si
      $\rg(A)=n$.
  \end{enumerate}
\end{theorem}

\begin{proof}
  \emph{(i)} Le système $A\vec{x}=\vec{b}$ est compatible si et seulement
  si $\vec{b}$ appartient à l'espace des colonnes de $A$, c'est-à-dire si
  $\vec{b}\in\im(A)$. Or $\vec{b}\in\im(A)$ si et seulement si l'ajout de
  $\vec{b}$ comme colonne supplémentaire ne change pas la dimension de
  l'espace des colonnes, c'est-à-dire
  $\rg(A\mid\vec{b})=\rg(A)$.

  Inversement, si $\vec{b}\notin\im(A)$, alors $\vec{b}$ est linéairement
  indépendant de l'espace des colonnes de $A$, et
  $\rg(A\mid\vec{b})=\rg(A)+1>\rg(A)$.

  \emph{(ii)} Posons $r=\rg(A)$ et considérons l'application linéaire
  $f\colon\K^n\to\K^m$, $\vec{x}\mapsto A\vec{x}$. On a $\im(f)$
  de dimension $r$, et le théorème du rang
  (cf.~
ef{ch:applications-lineaires}) donne~:
  \[
    \dim(\Ker(f))=n-\dim(\im(f))=n-r.
  \]
  Or $\Ker(f)$ est exactement l'ensemble des solutions du système homogène.

  \emph{(iii)} La solution est unique si et seulement si l'ensemble des
  solutions du système homogène est réduit à $\set{\vzero}$, c'est-à-dire
  $\dim(\Ker(f))=0$, ce qui équivaut à $r=n$.
\end{proof}

\begin{corollary}{Nombre de paramètres libres}{cor:params-libres}
  Si le système $A\vec{x}=\vec{b}$ ($A\in\Mat_{m,n}(\K)$) est compatible,
  l'ensemble des solutions est paramétré par $n-\rg(A)$ paramètres libres.
\end{corollary}

% ============================================================
\section{Structure de l'ensemble des solutions}\label{sec:structure-solutions}
% ============================================================

\begin{theorem}{Structure affine des solutions}{thm:structure-solutions}
  Soit $A\vec{x}=\vec{b}$ un système linéaire compatible, et soit
  $\vec{x}_p$ une solution particulière. Alors l'ensemble des solutions
  est le sous-espace affine\index{sous-espace affine}~:
  \begin{equation}\label{eq:structure-sol}
    \cS = \vec{x}_p + \Ker(A) = \setcond{\vec{x}_p+\vec{h}}{\vec{h}\in\Ker(A)}.
  \end{equation}
  Autrement dit, toute solution $\vec{x}$ du système s'écrit de manière
  unique~:
  \[
    \vec{x} = \vec{x}_p + \vec{x}_h,
  \]
  où $\vec{x}_p$ est une solution particulière de $A\vec{x}=\vec{b}$ et
  $\vec{x}_h$ est une solution du système homogène $A\vec{x}=\vzero$.
\end{theorem}

\begin{proof}
  Montrons la double inclusion.

  \emph{($\supseteq$)} Soit $\vec{h}\in\Ker(A)$. Alors
  $A(\vec{x}_p+\vec{h})=A\vec{x}_p+A\vec{h}=\vec{b}+\vzero=\vec{b}$,
  donc $\vec{x}_p+\vec{h}\in\cS$.

  \emph{($\subseteq$)} Soit $\vec{x}\in\cS$. Posons
  $\vec{h}=\vec{x}-\vec{x}_p$. Alors
  $A\vec{h}=A\vec{x}-A\vec{x}_p=\vec{b}-\vec{b}=\vzero$, donc
  $\vec{h}\in\Ker(A)$, et $\vec{x}=\vec{x}_p+\vec{h}\in\vec{x}_p+\Ker(A)$.

  \emph{Unicité de la décomposition.} Si $\vec{x}=\vec{x}_p+\vec{h}_1=
  \vec{x}_p+\vec{h}_2$ avec $\vec{h}_1,\vec{h}_2\in\Ker(A)$, alors
  $\vec{h}_1=\vec{h}_2$.
\end{proof}

\begin{remark}{Interprétation géométrique}{rem:interpretation-geom-sol}
  Le 
ef{thm:thm:structure-solutions} dit que l'ensemble des solutions
  est une \og copie translatée\fg{} du noyau de $A$. En dimension finie~:
  \begin{itemize}
    \item Si $\rg(A)=n$~: $\cS=\set{\vec{x}_p}$ est un point.
    \item Si $\rg(A)=n-1$~: $\cS$ est une droite affine.
    \item Si $\rg(A)=n-2$~: $\cS$ est un plan affine.
    \item En général~: $\cS$ est un sous-espace affine de dimension
      $n-\rg(A)$.
  \end{itemize}
\end{remark}

\begin{figure}[ht]
\centering
\begin{tikzpicture}[scale=0.9, >=Stealth, 3d view={115}{25}]
  % Plan homogène passant par l'origine
  \filldraw[blue!15, opacity=0.6]
    (-2,-2,0) -- (2,-2,0) -- (2,2,0) -- (-2,2,0) -- cycle;
  \node[blue!70!black, font=\small] at (2.3,1.5,0) {$\Ker(A)$};

  % Plan affine (translaté)
  \filldraw[red!15, opacity=0.6]
    (-2,-2,2) -- (2,-2,2) -- (2,2,2) -- (-2,2,2) -- cycle;
  \node[red!70!black, font=\small] at (2.3,1.5,2) {$\vec{x}_p+\Ker(A)$};

  % Vecteur xp
  \draw[->, thick, black!60!green] (0,0,0) -- (0,0,2)
    node[midway, right, font=\small] {$\vec{x}_p$};

  % Origine
  \filldraw[black] (0,0,0) circle (1.5pt) node[below left, font=\small] {$\vzero$};
  \filldraw[black!60!green] (0,0,2) circle (1.5pt) node[above left, font=\small] {$\vec{x}_p$};
\end{tikzpicture}
\caption{L'ensemble des solutions $\cS=\vec{x}_p+\Ker(A)$ est un translaté du noyau.}
\label{fig:structure-solutions}
\end{figure}

% ============================================================
\section{Représentation paramétrique des solutions}\label{sec:parametrique}
% ============================================================

Lorsqu'un système compatible admet des variables libres, on exprime les
\emph{variables liées}\index{variable!liée} (correspondant aux colonnes
pivots) en fonction des \emph{variables libres}\index{variable!libre}
(correspondant aux colonnes non pivots).

\begin{example}{Représentation paramétrique}{ex:param-sol}
  Considérons le système de rang $2$ à $4$ inconnues~:
  \[
    \left(\begin{array}{cccc|c}
      1 & 2 & 0 & -1 & 3\\
      0 & 0 & 1 & 3 & -2
    \end{array}\right).
  \]
  Les colonnes pivots sont les colonnes $1$ et $3$~; les variables libres
  sont $x_2$ et $x_4$. En posant $x_2=s$, $x_4=t$~:
  \[
    \vec{x}=\begin{pmatrix}3-2s+t\\s\\-2-3t\\t\end{pmatrix}
    =\underbrace{\begin{pmatrix}3\\0\\-2\\0\end{pmatrix}}_{\vec{x}_p}
    +s\underbrace{\begin{pmatrix}-2\\1\\0\\0\end{pmatrix}}_{\vec{v}_1}
    +t\underbrace{\begin{pmatrix}1\\0\\-3\\1\end{pmatrix}}_{\vec{v}_2},
    \qquad s,t\in\K.
  \]
  Les vecteurs $\vec{v}_1$ et $\vec{v}_2$ forment une base de $\Ker(A)$,
  et $\cS$ est un plan affine de $\K^4$.
\end{example}

% ============================================================
\section{Interprétation géométrique}\label{sec:interpretation-geom}
% ============================================================

\subsection{Systèmes \texorpdfstring{$2\times 2$}{2x2}}\label{subsec:geom-2x2}

Chaque équation linéaire en deux inconnues $a_1 x_1+a_2 x_2=b$ définit
une droite\index{droite} dans $\R^2$. Un système de deux équations
correspond à l'intersection de deux droites~:
\begin{itemize}
  \item \textbf{Un point} si les droites sont sécantes ($\rg(A)=2$).
  \item \textbf{Une droite} si les droites sont confondues ($\rg(A)=1$
    et $\rg(A\mid\vec{b})=1$).
  \item \textbf{L'ensemble vide} si les droites sont parallèles
    distinctes ($\rg(A)=1$ et $\rg(A\mid\vec{b})=2$).
\end{itemize}

\subsection{Systèmes \texorpdfstring{$3\times 3$}{3x3}}\label{subsec:geom-3x3}

Chaque équation linéaire en trois inconnues définit un plan\index{plan}
dans $\R^3$. Un système de trois équations correspond à l'intersection de
trois plans.

\begin{figure}[ht]
\centering
\begin{tikzpicture}[scale=1.0, >=Stealth, 3d view={115}{20}]
  % Plan 1
  \filldraw[blue!20, opacity=0.5]
    (-2,-2,{-2*0.5+(-2)*0.3+1}) -- (2,-2,{2*0.5+(-2)*0.3+1})
    -- (2,2,{2*0.5+2*0.3+1}) -- (-2,2,{-2*0.5+2*0.3+1}) -- cycle;
  \node[blue!70!black, font=\footnotesize] at (2.3,2,{2*0.5+2*0.3+1}) {$\pi_1$};
  % Plan 2
  \filldraw[red!20, opacity=0.5]
    (-2,-2,{-2*(-0.3)+(-2)*0.7+0.5}) -- (2,-2,{2*(-0.3)+(-2)*0.7+0.5})
    -- (2,2,{2*(-0.3)+2*0.7+0.5}) -- (-2,2,{-2*(-0.3)+2*0.7+0.5}) -- cycle;
  \node[red!70!black, font=\footnotesize] at (2.3,2,{2*(-0.3)+2*0.7+0.5}) {$\pi_2$};
  % Plan 3
  \filldraw[green!20, opacity=0.5]
    (-2,-2,{-2*0.2+(-2)*(-0.4)+2}) -- (2,-2,{2*0.2+(-2)*(-0.4)+2})
    -- (2,2,{2*0.2+2*(-0.4)+2}) -- (-2,2,{-2*0.2+2*(-0.4)+2}) -- cycle;
  \node[green!60!black, font=\footnotesize] at (2.3,2,{2*0.2+2*(-0.4)+2}) {$\pi_3$};
  % Point d'intersection (approximatif)
  \filldraw[black] (0,0,1.2) circle (2pt)
    node[above right, font=\small] {solution unique};
\end{tikzpicture}
\caption{Trois plans en position générale dans $\R^3$ se coupent en un unique point.}
\label{fig:3plans}
\end{figure}

% ============================================================
\section{Applications}\label{sec:applications-systemes}
% ============================================================

\subsection{Réseaux de flux}\label{subsec:reseaux-flux}

\begin{example}{Réseau routier}{ex:reseau-routier}
  Considérons un réseau de quatre carrefours $A$, $B$, $C$, $D$ reliés par
  des routes à sens unique. À chaque carrefour, le flux entrant doit égaler
  le flux sortant (conservation du flux)\index{conservation du flux}.

  \begin{center}
  \begin{tikzpicture}[>=Stealth, node distance=2.5cm]
    \node[draw, circle] (A) {$A$};
    \node[draw, circle, right of=A] (B) {$B$};
    \node[draw, circle, below of=B] (C) {$C$};
    \node[draw, circle, below of=A] (D) {$D$};
    \draw[->] (A) -- node[above, font=\small] {$x_1$} (B);
    \draw[->] (B) -- node[right, font=\small] {$x_2$} (C);
    \draw[->] (C) -- node[below, font=\small] {$x_3$} (D);
    \draw[->] (D) -- node[left, font=\small] {$x_4$} (A);
    % Flux externes
    \draw[->] (-1.5,0) -- node[above, font=\small] {$100$} (A);
    \draw[->] (B) -- node[above, font=\small] {$60$} ++(1.5,0);
    \draw[->] ++($(C)+(1.5,0)$) -- node[above, font=\small] {$50$} (C);
    \draw[->] (D) -- node[below, font=\small] {$90$} ++(-1.5,0);
  \end{tikzpicture}
  \end{center}

  Les équations de conservation sont~:
  \[
    \left\{
    \begin{aligned}
      A&: & 100+x_4 &= x_1, \\
      B&: & x_1 &= x_2+60, \\
      C&: & x_2+50 &= x_3, \\
      D&: & x_3 &= x_4+90.
    \end{aligned}
    \right.
    \quad\Longleftrightarrow\quad
    \left\{
    \begin{aligned}
      x_1 - x_4 &= 100, \\
      x_1 - x_2 &= 60, \\
      x_2 - x_3 &= -50, \\
      x_3 - x_4 &= 90.
    \end{aligned}
    \right.
  \]
  La matrice augmentée se réduit en~:
  \[
    \left(\begin{array}{cccc|c}
      1 & 0 & 0 & -1 & 100\\
      0 & 1 & 0 & -1 & 40\\
      0 & 0 & 1 & -1 & 90\\
      0 & 0 & 0 & 0 & 0
    \end{array}\right).
  \]
  Le rang est $3$ avec $4$ inconnues~: il y a un paramètre libre $x_4=t$.
  L'ensemble des solutions est~:
  \[
    (x_1,x_2,x_3,x_4)=(100+t,\;40+t,\;90+t,\;t),\qquad t\geqslant 0.
  \]
  La contrainte $t\geqslant 0$ (flux positifs) n'est pas une condition
  linéaire, mais elle limite le paramètre.
\end{example}

\subsection{Circuits électriques}\label{subsec:circuits}

\begin{example}{Lois de Kirchhoff}{ex:kirchhoff}
  Dans un circuit à deux mailles partageant une branche commune, les lois
  de Kirchhoff (loi des n\oe{}uds et loi des mailles) donnent un système
  linéaire dont les inconnues sont les courants $I_1$, $I_2$, $I_3$.
  Par exemple, avec des résistances $R_1=2\,\Omega$, $R_2=4\,\Omega$,
  $R_3=6\,\Omega$ et des tensions $V_1=12\,$V, $V_2=8\,$V~:
  \[
    \left\{
    \begin{aligned}
      I_1 - I_2 - I_3 &= 0 \quad\text{(loi des n\oe{}uds)},\\
      2I_1 + 6I_3 &= 12 \quad\text{(maille 1)},\\
      4I_2 - 6I_3 &= 8 \quad\text{(maille 2)}.
    \end{aligned}
    \right.
  \]
  La résolution par Gauss donne
  $(I_1, I_2, I_3)=\bigl(\frac{18}{5},\,\frac{13}{5},\,\frac{1}{1}\bigr)
  =\bigl(3{,}6\;;\;2{,}6\;;\;1{,}0\bigr)$ A.
\end{example}

\subsection{Aperçu~: moindres carrés}\label{subsec:moindres-carres}

Lorsqu'un système surdéterminé $A\vec{x}=\vec{b}$ (avec $m>n$) n'a pas de
solution exacte, on cherche le vecteur $\vec{x}^{*}$ qui minimise
$\norm{A\vec{x}-\vec{b}}^2$. Ce problème de \emph{moindres carrés}%
\index{moindres carrés} conduit aux \emph{équations normales}\index{équations normales}~:
\begin{equation}\label{eq:equations-normales}
  \transpose{A}A\,\vec{x}^{*}=\transpose{A}\vec{b},
\end{equation}
qui forment un système linéaire carré de taille $n\times n$. Ce sujet sera
développé dans le chapitre sur les espaces euclidiens.

% ============================================================
\section*{Résumé du chapitre}
\addcontentsline{toc}{section}{Résumé du chapitre}
% ============================================================

\begin{enumerate}[label=\textbf{\arabic*.}]
  \item Un \textbf{système linéaire}\index{système linéaire} $A\vec{x}=\vec{b}$
    se représente par sa matrice augmentée $(A\mid\vec{b})$.
  \item Les \textbf{opérations élémentaires} sur les lignes (échange,
    multiplication par un scalaire non nul, ajout d'un multiple) ne
    changent pas l'ensemble des solutions.
  \item L'\textbf{élimination de Gauss} transforme la matrice augmentée en
    \textbf{forme échelonnée}~; le prolongement de Gauss--Jordan produit la
    \textbf{forme échelonnée réduite} (unique).
  \item Le \textbf{rang} $\rg(A)$ est le nombre de pivots dans une forme
    échelonnée.
  \item \textbf{Théorème de Rouché--Capelli~:} le système est compatible si
    et seulement si $\rg(A)=\rg(A\mid\vec{b})$~; dans ce cas, l'ensemble
    des solutions est paramétré par $n-\rg(A)$ variables libres.
  \item \textbf{Structure des solutions~:} si $\vec{x}_p$ est une solution
    particulière, alors $\cS=\vec{x}_p+\Ker(A)$ (sous-espace affine).
  \item Les systèmes linéaires modélisent des problèmes concrets~: réseaux
    de flux, circuits électriques, moindres carrés, etc.
\end{enumerate}

% ============================================================
\section*{Exercices}
\addcontentsline{toc}{section}{Exercices}
% ============================================================

\begin{exercise}{Résolution par Gauss \basic}{exo:gauss-basique}
  Résoudre le système suivant par élimination de Gauss~:
  \[
    \left\{
    \begin{aligned}
      x_1 + 2x_2 + 3x_3 &= 9,\\
      2x_1 - x_2 + x_3 &= 8,\\
      3x_1 + x_2 - 2x_3 &= 3.
    \end{aligned}
    \right.
  \]
\end{exercise}

\begin{exercise}{Forme échelonnée réduite \basic}{exo:rref-basique}
  Mettre la matrice suivante sous forme échelonnée réduite~:
  \[
    A=\begin{pmatrix}
      1 & 3 & 5 & 7\\
      2 & 4 & 6 & 8\\
      3 & 5 & 7 & 9
    \end{pmatrix}.
  \]
  En déduire le rang de $A$.
\end{exercise}

\begin{exercise}{Système homogène \basic}{exo:homogene-basique}
  Déterminer l'ensemble des solutions du système homogène~:
  \[
    \left\{
    \begin{aligned}
      x_1 - x_2 + 2x_3 - x_4 &= 0,\\
      2x_1 + x_2 - x_3 + x_4 &= 0,\\
      x_1 + 2x_2 - 3x_3 + 2x_4 &= 0.
    \end{aligned}
    \right.
  \]
  Donner une base du noyau de la matrice des coefficients.
\end{exercise}

\begin{exercise}{Compatibilité d'un système \basic}{exo:compatibilite}
  Déterminer pour quelles valeurs du paramètre $a\in\R$ le système
  suivant est compatible~:
  \[
    \left\{
    \begin{aligned}
      x_1 + x_2 + x_3 &= 1,\\
      x_1 + 2x_2 + 3x_3 &= 3,\\
      x_1 + 3x_2 + ax_3 &= 5.
    \end{aligned}
    \right.
  \]
\end{exercise}

\begin{exercise}{Système paramétrique \intermediate}{exo:param-systeme}
  Résoudre en fonction du paramètre $\lambda\in\R$~:
  \[
    \left\{
    \begin{aligned}
      \lambda x_1 + x_2 + x_3 &= 1,\\
      x_1 + \lambda x_2 + x_3 &= 1,\\
      x_1 + x_2 + \lambda x_3 &= 1.
    \end{aligned}
    \right.
  \]
  Distinguer les cas selon les valeurs de $\lambda$.
\end{exercise}

\begin{exercise}{Calcul de rang \intermediate}{exo:rang-calcul}
  Calculer le rang des matrices suivantes en fonction de $\alpha\in\R$~:
  \[
    A_1=\begin{pmatrix}
      1 & 2 & 3\\
      4 & 5 & 6\\
      7 & 8 & 9
    \end{pmatrix},\qquad
    A_2=\begin{pmatrix}
      1 & 1 & 1 & 1\\
      1 & 2 & 3 & 4\\
      1 & 3 & \alpha & 10
    \end{pmatrix}.
  \]
\end{exercise}

\begin{exercise}{Structure affine \intermediate}{exo:structure-affine}
  On considère le système $A\vec{x}=\vec{b}$ avec~:
  \[
    A=\begin{pmatrix}1&1&0&1\\0&1&1&1\\1&2&1&2\end{pmatrix},\qquad
    \vec{b}=\begin{pmatrix}2\\1\\3\end{pmatrix}.
  \]
  \begin{enumerate}[label=(\alph*)]
    \item Vérifier que le système est compatible.
    \item Trouver une solution particulière.
    \item Déterminer une base de $\Ker(A)$.
    \item Écrire l'ensemble des solutions sous la forme
      $\cS=\vec{x}_p+\Vect(\ldots)$.
  \end{enumerate}
\end{exercise}

\begin{exercise}{Réseau de flux \intermediate}{exo:reseau-flux}
  Dans le réseau ci-dessous, les flux entrants et sortants aux n\oe{}uds
  sont indiqués~:
  \[
    \begin{array}{ccc}
      & \xrightarrow{x_1} & \\[-3pt]
      A & & B\\[-3pt]
      & & \downarrow\, x_2\\[-3pt]
      & & C\\[-3pt]
      & \xleftarrow{x_3} &
    \end{array}
  \]
  Flux externe entrant en $A$~: $200$~; flux externe sortant en $B$~: $80$~;
  flux externe sortant en $C$~: $120$.
  \begin{enumerate}[label=(\alph*)]
    \item Écrire le système de conservation du flux.
    \item Résoudre le système et interpréter.
  \end{enumerate}
\end{exercise}

\begin{exercise}{Inversibilité et rang \intermediate}{exo:inversibilite-rang}
  Soit $A\in\Mat_n(\K)$ une matrice carrée. Montrer que les assertions
  suivantes sont équivalentes~:
  \begin{enumerate}[label=(\roman*)]
    \item $A$ est inversible.
    \item $\rg(A)=n$.
    \item Le système $A\vec{x}=\vec{b}$ a une unique solution pour tout
      $\vec{b}\in\K^n$.
    \item Le système homogène $A\vec{x}=\vzero$ n'a que la solution triviale.
  \end{enumerate}
\end{exercise}

\begin{exercise}{Gauss--Jordan et inverse \challenging}{exo:gauss-jordan-inverse}
  Utiliser l'élimination de Gauss--Jordan pour calculer l'inverse de la
  matrice~:
  \[
    A=\begin{pmatrix}
      1 & 2 & 1\\
      2 & 5 & 3\\
      1 & 3 & 3
    \end{pmatrix}.
  \]
  \textit{Indication~: appliquer Gauss--Jordan à la matrice augmentée
  $(A\mid I_3)$.}
\end{exercise}

\begin{exercise}{Systèmes et applications linéaires \challenging}{exo:syst-applin}
  Soit $f\colon\R^4\to\R^3$ l'application linéaire définie par~:
  \[
    f(x_1,x_2,x_3,x_4)=(x_1+x_2+x_3+x_4,\;
    2x_1+x_2-x_3+2x_4,\;
    3x_1+2x_2+x_4).
  \]
  \begin{enumerate}[label=(\alph*)]
    \item Écrire la matrice de $f$ dans les bases canoniques.
    \item Calculer $\rg(f)$ et $\dim(\Ker(f))$.
    \item Déterminer une base de $\Ker(f)$ et une base de $\im(f)$.
    \item L'équation $f(\vec{x})=(1,2,3)$ a-t-elle des solutions~? Si oui,
      les déterminer toutes.
  \end{enumerate}
\end{exercise}

\begin{exercise}{Théorème de Rouché--Capelli \challenging}{exo:rouche-capelli-preuve}
  Soit $A\in\Mat_{m,n}(\K)$ et $\vec{b}\in\K^m$. On suppose que
  $\rg(A)=\rg(A\mid\vec{b})=r$.
  \begin{enumerate}[label=(\alph*)]
    \item Montrer que si $r=n$, la solution est unique.
    \item Montrer que si $r<n$ et si $\vec{v}_1,\ldots,\vec{v}_{n-r}$ est
      une base de $\Ker(A)$, alors toute solution s'écrit~:
      \[
        \vec{x}=\vec{x}_p+t_1\vec{v}_1+\cdots+t_{n-r}\vec{v}_{n-r},
        \qquad t_1,\ldots,t_{n-r}\in\K.
      \]
    \item En déduire que le nombre de solutions est soit $0$, soit $1$,
      soit infini (lorsque $\K$ est infini).
  \end{enumerate}
\end{exercise}

\begin{exercise}{Système surdéterminé et moindres carrés \challenging}{exo:moindres-carres}
  On mesure les points $(0,1)$, $(1,2)$, $(2,4)$, $(3,5)$ et on cherche la
  droite $y=a+bx$ au sens des moindres carrés.
  \begin{enumerate}[label=(\alph*)]
    \item Écrire le système surdéterminé $A\vec{x}=\vec{b}$ correspondant.
    \item Former les équations normales $\transpose{A}A\,\vec{x}^{*}=\transpose{A}\vec{b}$.
    \item Résoudre et déterminer la droite de régression.
  \end{enumerate}
\end{exercise}

\printindex
